\section{Термодинамика}

\subsection{Температура}

\textbf{Температура ($T$)} --- макроскопический параметр, являющийся мерой средней кинетической энергии теплового движения молекул.

Метод измерения температуры основан на использовании термометрического свойства жидкости --- \textit{теплового расширения}. Предполагается, что объем жидкости $V$ изменяется линейно с температурой.

Для построения шкалы Цельсия выбираются две опорные точки при нормальном атмосферном давлении:
\begin{itemize}
    \item \textbf{0$^\circ$C}: Температура таяния льда. Измеряется объем жидкости $V_0$.
    \item \textbf{100$^\circ$C}: Температура кипения воды. Измеряется объем жидкости $V_{100}$.
\end{itemize}

Пусть при некоторой неизвестной температуре $t_x$ объем жидкости принял значение $V_x$. Исходя из предположения о линейности шкалы, искомая температура вычисляется по формуле:

\[
t_x = \frac{V_x - V_0}{V_{100} - V_0} \times 100^\circ\text{C}
\]

Если вместо объема мы измеряем высоту столба жидкости $h$ в трубке постоянного сечения ($V = S \cdot h$), формула принимает вид:
\[
t_x = \frac{h_x - h_0}{h_{100} - h_0} \times 100^\circ\text{C}
\]

\subsection{Распределение Максвелла}
Вероятность распределения скоростей молекул при температуре $T$:
\[
f(v) = 4\pi \left( \frac{m}{2\pi kT} \right)^{3/2} v^2 e^{-\frac{mv^2}{2kT}}
\]

Ниже представлен график функции распределения $f(v)$ с указанием характерных скоростей:

\begin{center}
\begin{tikzpicture}
\begin{axis}[
    height=6cm, width=10cm,
    axis lines = middle,
    axis line style={thick, black, dashed},
    xlabel = {$v$},
    ylabel = {$f(v)$},
    xlabel style={right},
    ylabel style={above},
    domain = 0:3.5,
    samples = 100,
    ytick = \empty,
    xtick = \empty,
    enlarge y limits={upper, value=0.2},
    clip=false
]
    % Формула функции распределения
    \addplot [very thick, black] {4/sqrt(pi) * x^2 * exp(-x^2)};

    % 1. Мода (Наиболее вероятная скорость) v_p
    \draw [dashed, black, thick] (axis cs:1, 0) -- (axis cs:1, 0.83);
    \node [anchor=north, black] at (axis cs:1, 0) {$v_p$};

    % 2. Средняя скорость <v>
    \draw [dashed, black, thick] (axis cs:1.2, 0) -- (axis cs:1.2, 0.77);
    \node [anchor=south, black] at (axis cs:1.2, 0.8) {$\langle v \rangle$};

    % 3. Среднеквадратичная скорость v_rms
    \draw [dashed, black, thick] (axis cs:1.4, 0) -- (axis cs:1.4, 0.62);
    \node [anchor=north, black] at (axis cs:1.45, 0) {$v_{\sigma}$};

\end{axis}
\end{tikzpicture}
\end{center}

\begin{center}
\begin{tblr}{
  colspec = {X[c]X[l]},
  vlines,
  hlines,
  row{1} = {font=\bfseries,mode=text},
  row{odd} = {bg=gray!10},
  cells = {m},
  rowsep = 6pt
}
\textbf{Формула} & \textbf{Название} \\
$v_p= \sqrt{\dfrac{2kT}{m}}$ & Мода \\
$\langle v \rangle = \sqrt{\dfrac{8kT}{\pi m}}$ & Средняя скорость \\
$v_{\sigma} = \sqrt{\langle v^2 \rangle} = \sqrt{\dfrac{3kT}{m}}$ & Среднеквадратичная скорость \\
\end{tblr}
\end{center}

\subsection{Теплообмен}

Перенос внутренней энергии между телами или частями системы осуществляется тремя основными способами: теплопроводностью, конвекцией и излучением.

\textbf{Теплопроводность} --- это перенос внутренней энергии за счёт неупорядоченного движения молекул, атомов, электронов при их непосредственном контакте, без переноса вещества.
\begin{itemize}
    \item \textit{В твёрдых телах:} основной механизм --- колебания кристаллической решётки и движение свободных электронов (в металлах).
    \item \textit{В жидкостях и газах:} передача энергии при столкновениях молекул.
\end{itemize}

Количественно описывается \textit{законом Фурье}:
\[
\diff Q = -\lambda\diff A\diff\tau\cdot\grad T
\]

\textbf{Тепловой поток ($\Phi$)} --- количество теплоты, проходящее через выбранную поверхность в единицу времени: 
\[
\Phi = \frac{dQ}{\diff\tau}=-\lambda\diff A\cdot\grad T
\]

\textbf{Плотность теплового потока} ($q$) --- количество теплоты, проходящее через единичную площадку в единицу времени:
\[
q = \frac{\diff\Phi}{\diff A} = -\lambda\cdot\grad T
\]

На практике коэффициент $\lambda$ часто является функцией температуры \( f(T) \). В этом случае уравнения теплопроводности становятся нелинейными.

\textbf{Теплоёмкость ($C$)} --- характеристика вещества, равная отношению сообщённой теплоты к вызванному ею изменению температуры: 
\[
C = \frac{\delta Q}{dT}
\]

\textbf{Конвекция} --- перенос теплоты потоками жидкости или газа. Обязательно сопровождается переносом вещества:

\begin{itemize}
    \item \textbf{естественная} конвекция возникает из-за разницы плотностей в неоднородно нагретой среде \textit{(движение воздуха у батареи)}.
    \item \textbf{вынужденная} конвекция вызывается внешними силами \textit{(вентилятор)}.
\end{itemize}

\textbf{Тепловое излучение} --- перенос энергии электромагнитными волнами, испускаемыми нагретыми телами за счёт их внутренней энергии. Не требует среды для распространения (работает в вакууме).

\subsection{Идеальный газ}

\subsubsection{Определение}

\textbf{Идеальный газ} --- это теоретическая модель газа со свойствами:
\begin{itemize}
  \item Собственный объем молекул пренебрежимо мал по сравнению с объемом сосуда (молекулы --- материальные точки).
  \item Отсутствуют силы межмолекулярного взаимодействия на расстоянии.
  \item Молекулы сталкиваются абсолютно упруго \textit{(стенки, между собой)}.
\end{itemize}

\subsubsection{Уравнения}

\begin{center}
\begin{tblr}{
  colspec = {X[1,c]X[2,l]},
  vlines,
  hlines,
  row{1} = {font=\bfseries,mode=text},
  row{odd} = {bg=gray!10},
  cells = {m},
  rowsep = 6pt
}
\textbf{Уравнение} & \textbf{Название и условия} \\
$\displaystyle P V = \nu R T$ & Уравнение Менделеева--Клапейрона \\
$\displaystyle \frac{P V}{T} = \text{const}$ & Уравнение Клапейрона ($m = \text{const}$) \\
$P V = \text{const}$ & Закон Бойля--Мариотта ($T = \text{const}$) \\
$\displaystyle \frac{V}{T} = \text{const}$ & Закон Гей-Люссака ($P = \text{const}$) \\
$\displaystyle \frac{P}{T} = \text{const}$ & Закон Шарля ($V = \text{const}$) \\
\end{tblr}
\end{center}

\begin{itemize}
  \item $R = k_B N_A \approx 8.314 \ \text{Дж/(моль·К)}$ --- универсальная газовая постоянная
  \item $k_B = 1.38 \times 10^{-23}$ Дж/К --- постоянная Больцмана
\end{itemize}

\subsubsection{Закон Авогадро}

\textbf{Закон Авогадро} утверждает, что в равных объемах различных газов при одинаковых температуре и давлении содержится одинаковое число молекул.
\begin{itemize}
    \item \textbf{Молярный объем:} При нормальных условиях ($P_0 = 10^5$ Па, $T_0 = 0^\circ$C) объем одного моля любого газа составляет $V_m \approx 22.4$ л/моль.
    \item \textbf{Число Авогадро:} $N_A \approx 6.022 \times 10^{23}$ моль$^{-1}$.
\end{itemize}
Связь между числом частиц $N$ и количеством вещества $\nu$:
\[
N = \nu \cdot N_A = \frac{m}{M} N_A
\]